% SOURCE_STATEMENT: "Explain the mechanism producing limit cycles in shallow cumulus organization. Determine how external forcings affect their size, frequency, and amplitude; whether they respond quasi-statically to time-dependent forcing; how the diurnal radiative cycle and realistic open boundaries modify them; and how relevant they are to cloud feedback."
% Preparation and final self-audit: CONTRIBUTING.md prompts C and D; source/audit artifacts retained in the local physics_statements_2026_10_11 batch.
% !TeX program = lualatex
\documentclass[11pt]{article}
\usepackage[margin=25mm]{geometry}
\usepackage{fontspec}
\setmainfont{Latin Modern Roman}
\usepackage{amsmath,amssymb,mathtools}
\usepackage{unicode-math}
\setmathfont{Latin Modern Math}
\usepackage{xurl,hyperref}
\hypersetup{colorlinks=true,urlcolor=blue}
\setcounter{secnumdepth}{0}
\setlength{\parindent}{0pt}
\setlength{\parskip}{5pt}
\setlength{\emergencystretch}{3em}
\newcommand{\sref}[2]{\hyperlink{source:#1:#2}{[S#2]}}
\newcommand{\cataloguescope}[1]{\paragraph{Recorded target.}#1}
\newcommand{\R}{\mathbb R}
\newcommand{\C}{\mathbb C}
\newcommand{\N}{\mathbb N}
\newcommand{\Z}{\mathbb Z}
\newcommand{\Q}{\mathbb Q}
\newcommand{\F}{\mathbb F}
\newcommand{\eps}{\varepsilon}
\begin{document}
% UnsolvedMath ID 30005988; source catalogue code OWR-14298586-004
\setcounter{section}{30005987}
\section{Cloud-Organization Limit Cycles}\label{30005988}
{\small UnsolvedMath ID 30005988 · Mathematical Physics}
\subsection{Definitions and mathematical statement}

\paragraph{Shared notation.}$\mathbb N=\{1,2,\ldots\}$; the standard number fields have their usual meanings. All additional conventions are stated locally.


A limit cycle is an isolated periodic orbit of a dynamical model of cloud organization. Its size, period and amplitude describe different observables. Quasistatic response means tracking the family of equilibrated or periodic states under slow forcing. Diurnal radiation and open boundaries alter the forcing and exchange of mass and energy.

\paragraph{Catalogue wording retained for traceability.}
Explain the mechanism producing limit cycles in shallow cumulus organization. Determine how external forcings affect their size, frequency, and amplitude; whether they respond quasi-statically to time-dependent forcing; how the diurnal radiative cycle and realistic open boundaries modify them; and how relevant they are to cloud feedback.

\paragraph{Scope and source qualification.}All forcing-response questions are retained, including whether time-dependent forcing is followed quasistatically.
\subsection{Short English statement}
Explain cloud-organization cycles and how forcing, daily radiation and open boundaries change their scales and climate feedback.
\subsection{Sources}
\begin{itemize}
\item[\hypertarget{source:30005988:1}{S1}] ULAM.AI and the UnsolvedMath Contributors, UnsolvedMath, record 30005988 (OWR-14298586-004). Catalogue curation under CC BY 4.0; definitions and English exposition added locally.
\url{https://huggingface.co/datasets/ulamai/UnsolvedMath}.
\item[\hypertarget{source:30005988:2}{S2}] Model Hierarchies in Atmosphere, Ocean, and Climate Sciences (2024), Oberwolfach Reports, DOI 10.4171/owr/2024/31, pp. 1731--1786. Original report and mathematical context. Incorporated report material remains under the source CC BY-SA 4.0 terms unless otherwise noted.
\url{https://doi.org/10.4171/owr/2024/31}.
\end{itemize}
\label{end:30005988}
\end{document}
